% BHU PYQ -- Professional Exam Template
\documentclass[11pt,a4paper]{article}

\usepackage[a4paper,top=2.5cm,bottom=2.5cm,left=2.8cm,right=2.8cm,headheight=14pt]{geometry}
\usepackage{amsmath,amssymb}
\usepackage[T1]{fontenc}
\usepackage[utf8]{inputenc}
\usepackage{lmodern}
\usepackage{microtype}
\usepackage{fancyhdr}
\usepackage{enumitem}
\usepackage{listings}
\usepackage{hyperref}
\usepackage{lastpage}
\usepackage{parskip}

\hypersetup{colorlinks=true,urlcolor=black,linkcolor=black,
  pdftitle={BVCA-301: Analog Digital Communication -- BHU PYQ}}

\pagestyle{fancy}
\fancyhf{}
\renewcommand{\headrulewidth}{0.4pt}
\renewcommand{\footrulewidth}{0.4pt}
\fancyhead[L]{\small\textit{Banaras Hindu University}}
\fancyhead[R]{\small\textit{BVCA-301: Analog Digital Communication}}
\fancyfoot[C]{\small Page \thepage\ of \pageref{LastPage}}

\lstset{
  basicstyle=\small\ttfamily,
  frame=single,
  breaklines=true,
  columns=flexible,
  keepspaces=true
}

\newlist{parts}{enumerate}{2}
\setlist[parts,1]{label=(\alph*),leftmargin=*,itemsep=4pt,topsep=3pt}
\setlist[parts,2]{label=(\roman*),leftmargin=*,itemsep=2pt,topsep=2pt}
\newcommand{\pts}[1]{\hfill\textit{[#1]}}

\begin{document}

\begin{center}
  {\large\bfseries BANARAS HINDU UNIVERSITY}\\[2pt]
  {\normalsize B.Voc. Semester III Examination, 2023-24}\\[6pt]
  \rule{0.8\linewidth}{0.5pt}\\[6pt]
  {\large\bfseries Paper: BVCA-301 --- Analog Digital Communication}\\[4pt]
  {\normalsize Time: Three hours \hfill Maximum Marks: 70}
\end{center}

\rule{\linewidth}{0.4pt}

\smallskip
\noindent\textit{Instructions: Answer the questions in each section as instructed.}

\bigskip

BVCA-301 : Analog and Digital Communication
Note: Answer the questions in each section as instructed.
All questions carry marks as indicated. |

\bigskip\noindent{\large\bfseries SECTION --- A}
\rule{\linewidth}{0.4pt}\smallskip

Note: Long-Answer Questions each in 500 words approximately (Answer any 2 out .
of 4; Each question carries 12 marks). 2x12=24 |

\medskip\noindent\textbf{Question 1.} Define amplitude modulation and describe it in detail with suitable diagrams. 3

\medskip\noindent\textbf{Question 2.} Determine the following for DSBSC (i) Bandwidth (2Marks) (ii) Time domain representation .
(5Marks) (iii) Frequency spectrum Waveform (3Marks) (iv) Advantages, Disadvantages and |
Applications(2Marks)

\medskip\noindent\textbf{Question 3.} Fora modulation coefficient m = 0.2 and an unmodulated carrier power Pc = 1000W, Examine
the total sideband power, upper \& lower side band power, modulated carrier power and total
transmitted power. :

\medskip\noindent\textbf{Question 4.} State the principle of Angle Modulation. Derive phase deviation, modulation index, frequency
deviation and percent modulation. :

\bigskip\noindent{\large\bfseries SECTION --- B}
\rule{\linewidth}{0.4pt}\smallskip


\medskip\noindent\textbf{Question 1.} Short-Answer Questions each in 200 words approximately. Answer any 6 out
of 9; Each question carries 6 marks. 6x6=36
\begin{parts}
  \item Define sampling and express Nyquist sampling rate.
  \item What are the elementary signals define with graphs and mathematical functions?
  \item Define System and how to check the linearity of a system? Give suitable example with proof.
  \item What is Fourier Transform? Calculate the Fourier Transform of the following signals (i)
  e *u(t) where u(t) is the unit step signal (ii) e °u(t) + $e^{at}u(-t)$.
  (€) If x(t) = u(t) is an unit step signal then draw the following (i) 3x(t — 2) (ii) 5x(—2t +1)
  \item Classify the different methods of Pulse modulation techniques.
  \item Discuss the generation of DSB SC Waves. :
  \item A 1000kHz carrier is simultaneously modulated with 300Hz, 800Hz and 2kHz audio sine waves.
  Select the frequencies present in the output.
  \item Write the advantages of digital communication over analog communication.
\end{parts}

\bigskip\noindent{\large\bfseries SECTION --- C}
\rule{\linewidth}{0.4pt}\smallskip

Note: Very Short-Answer Questions (Answer all 10; Each question carries 1 marks).
1x10=10
) Analog signals cannot be sent using digital techniques. (True (T)/False (F))
  \begin{parts}
    \item Digitizing a signal can reduce distortion. (True (T)/False (F))
    \item Digitization removes noise and distortion from analog signals. (True (T)/False (F))
    \item A digital signal can be changed from a | to a 0 by noise. (True (T)/False (F :
    g 8 )
    \item Unlike analog, digital communications is not band-limited. (True (T)/False (F))
    \item — The amount of digital data that can be sent is limited by the channel capacity (C). (True (T)/False
    wea ()
    \item What are the causal and non-causal signals?
    \item What is the bandwidth of an AM signal?
    \item | Write the full form of DSB-FC, DSB-SC, and SSB-SC.
    \item Write the Nyquist criterion for sampling.
  \end{parts}

\vfill
\rule{\linewidth}{0.4pt}
\begin{center}\small
\href{https://bhu-exam.vercel.app/}{Click here to visit our website}\\[4pt]\href{https://me-aryan.vercel.app/}{Click here to visit our contributor}\\[4pt]\href{https://quashan-me.vercel.app/}{Click here to visit our contributor}
\end{center}

\end{document}
